\documentclass[UTF8,11pt,a4paper]{ctexart}

\usepackage[a4paper,margin=2.3cm]{geometry}
\usepackage{booktabs}
\usepackage{longtable}
\usepackage{array}
\usepackage{enumitem}
\usepackage{xcolor}
% Treat underscores in file names and identifiers as text characters instead
% of mathematical subscript operators.
\usepackage[strings]{underscore}
\usepackage{hyperref}
\usepackage{fancyhdr}
\usepackage{listings}

\hypersetup{
  colorlinks=true,
  linkcolor=blue!55!black,
  urlcolor=blue!65!black,
  pdftitle={vLLM Ascend 项目目录与模块说明},
  pdfauthor={项目代码分析整理}
}

\pagestyle{fancy}
\fancyhf{}
\fancyhead[L]{vLLM Ascend 项目目录与模块说明}
\fancyhead[R]{\thepage}
\setlength{\headheight}{14pt}

\setlist[itemize]{leftmargin=2em,itemsep=0.3em}
\setlist[enumerate]{leftmargin=2em,itemsep=0.3em}
\renewcommand{\arraystretch}{1.25}
% These commands are robust because \folder is also used in section titles.
\DeclareRobustCommand{\folder}[1]{\texttt{\detokenize{#1}}}
\DeclareRobustCommand{\code}[1]{\texttt{\detokenize{#1}}}
% PDF bookmarks are plain strings; strip the formatting commands there.
\pdfstringdefDisableCommands{%
  \def\folder#1{#1}%
  \def\code#1{#1}%
}

\definecolor{codebg}{RGB}{246,248,250}
\definecolor{codeframe}{RGB}{210,215,220}
\lstset{
  basicstyle=\ttfamily\small,
  backgroundcolor=\color{codebg},
  frame=single,
  rulecolor=\color{codeframe},
  columns=fullflexible,
  keepspaces=true,
  showstringspaces=false,
  breaklines=true,
  tabsize=2
}

\title{\bfseries vLLM Ascend 项目目录与模块说明}
\author{基于项目目录 \folder{C:/Users/Administrator/Desktop/找工作/vllm-ascend} 整理}
\date{\today}

\begin{document}

\maketitle

\begin{abstract}
本文档介绍 \folder{vllm-ascend} 项目的总体定位、顶层目录、核心 Python 包、底层 CANN/C++ 算子、测试和开发工具，并给出主执行流程以及 LoPT 并行 tokenization 功能可能涉及的代码位置。该项目不是独立的 HTTP 推理服务，而是上游 vLLM 的 Ascend NPU 硬件后端插件。
\end{abstract}

\tableofcontents
\newpage

\section{项目定位}

\folder{vllm-ascend} 的主要功能是把上游 vLLM 的模型推理框架适配到华为 Ascend NPU。上游 vLLM 负责 API、请求预处理、tokenization、通用调度和输出处理；本项目主要负责 NPU 平台注册、设备初始化、模型执行、Attention、KV Cache、采样、分布式通信和自定义算子。

项目可粗略划分为三层：

\begin{lstlisting}
工程、文档、测试和协作配置
            |
Python 运行时：vllm_ascend/
            |
C++ / AscendC / CANN 高性能算子：csrc/
\end{lstlisting}

真正影响在线推理运行的核心目录是 \folder{vllm_ascend/} 和 \folder{csrc/}。

\section{顶层目录}

\begin{longtable}{>{\raggedright\arraybackslash}p{0.21\textwidth} >{\raggedright\arraybackslash}p{0.58\textwidth} >{\centering\arraybackslash}p{0.12\textwidth}}
\toprule
\textbf{目录} & \textbf{主要作用} & \textbf{参与推理} \\
\midrule
\endfirsthead
\toprule
\textbf{目录} & \textbf{主要作用} & \textbf{参与推理} \\
\midrule
\endhead
\folder{.agents/} & Codex 或其他智能代理使用的项目技能、规则和说明。 & 否 \\
\folder{.claude/} & Claude 辅助开发相关说明。 & 否 \\
\folder{.gemini/} & Gemini 配置和代码风格说明。 & 否 \\
\folder{.git/} & Git 仓库元数据、对象和版本历史。 & 否 \\
\folder{.github/} & GitHub Actions、Issue 模板、公共 Action 和 CI 工作流。 & 否 \\
\folder{benchmarks/} & 吞吐、延迟、Serving 和算子性能测试及结果处理。 & 测试时 \\
\folder{cmake/} & 顶层 CMake 公共构建函数。 & 编译时 \\
\folder{csrc/} & C++、AscendC 和 CANN 自定义高性能算子。 & 是 \\
\folder{docs/} & 中英文文档、文档网站源码、主题覆盖和构建钩子。 & 否 \\
\folder{examples/} & 量化、PD 分离、动态负载均衡、强化学习等使用示例。 & 示例运行 \\
\folder{tests/} & 单元测试、端到端测试、夜间和每周测试。 & 测试时 \\
\folder{tools/} & Benchmark、请求发送、静态检查、回归定位和文档生成工具。 & 工具运行 \\
\folder{vllm_ascend/} & Ascend 插件的核心 Python 实现。 & 是 \\
\bottomrule
\end{longtable}

\section{核心 Python 包：\folder{vllm_ascend/}}

\subsection{\folder{worker/}：推理 Worker 与模型执行器}

该目录是运行时最核心的目录之一，负责把上游 Scheduler 产生的任务转换成 NPU 模型执行。

主要文件包括：

\begin{itemize}
  \item \folder{worker.py}：定义 \code{NPUWorker}，负责设备初始化、模型加载、KV Cache 初始化、推理执行和 Pipeline Parallel 通信。
  \item \folder{model_runner_v1.py}：当前主要模型执行路径，负责输入准备、Attention metadata、模型 forward、logits 和采样状态。
  \item \folder{npu_input_batch.py}：维护 NPU 上的请求批次和输入状态。
  \item \folder{block_table.py}：维护请求与 KV Cache block 的映射。
  \item \folder{pcp_utils.py}：Parallel Context Processing 相关逻辑。
  \item \folder{encoder_acl_graph.py}：Encoder 模型 ACL Graph 管理。
  \item \folder{v2/}：开发中的 v2 ModelRunner 实现。
\end{itemize}

其主要调用关系为：

\begin{lstlisting}
NPUWorker.execute_model()
        |
NPUModelRunner.execute_model()
        |
input_ids / positions / KV blocks
        |
Attention + model forward
        |
logits + token sampling
\end{lstlisting}

\subsection{\folder{attention/}：Ascend Attention 后端}

该目录实现不同模型和硬件场景下的 Attention 计算，包括：

\begin{itemize}
  \item \folder{attention_v1.py}：普通 Attention 后端。
  \item \folder{mla_v1.py}：Multi-head Latent Attention（MLA）。
  \item \folder{dsa_v1.py}：DSA Attention。
  \item \folder{sfa_v1.py}：SFA Attention。
  \item \folder{fa3_v1.py}：FA3 后端。
  \item \folder{attention_mask.py}：Attention mask 构造。
  \item \folder{context_parallel/}：上下文并行实现。
  \item \folder{kvcomp_attn/}：KV 压缩 Attention 相关实现。
\end{itemize}

这里处理的是 token ID 已经转为张量之后的 Attention 计算，不负责文本 tokenization。

\subsection{\folder{ops/}：Python 层 Ascend 算子}

该目录提供 Python 层的自定义算子、融合算子和算子注册，包括 Linear、RMSNorm、LayerNorm、RoPE、MoE、MLA、DSA、多模态 Encoder Attention、Vocabulary Parallel Embedding、权重预取及 Triton 算子。

\folder{register_custom_ops.py} 负责注册部分自定义算子。Python 层算子可能进一步调用 \folder{csrc/} 编译产生的底层扩展。

\subsection{\folder{core/}：Scheduler 与 KV Cache 扩展}

该目录扩展上游 vLLM 的调度和 KV Cache 管理能力，主要包含：

\begin{itemize}
  \item 动态 batch scheduler；
  \item 短请求优先 scheduler；
  \item recompute scheduler；
  \item profiling chunk scheduler；
  \item KV Cache 接口；
  \item 单类型 KV Cache manager。
\end{itemize}

它决定已经完成 tokenization 的请求如何进入模型，但不处理原始文本。

\subsection{\folder{compilation/}：图编译与 ACL Graph}

负责 Ascend 图编译、执行图管理和图优化，包括：

\begin{itemize}
  \item ACL Graph 捕获与重放；
  \item 编译器接口；
  \item Graph Fusion Pass；
  \item 编译 Pass 管理。
\end{itemize}

其目标是减少 Python 调度开销，并复用适合静态执行的 NPU 计算图。

\subsection{\folder{distributed/}：分布式推理与数据传输}

该目录负责 Ascend 分布式执行和跨进程数据传输：

\begin{itemize}
  \item \folder{device_communicators/}：设备通信实现；
  \item \folder{kv_transfer/}：KV Cache 传输、Prefill/Decode 分离、Mooncake 和 Ascend Store；
  \item \folder{weight_transfer/}：运行时权重传输；
  \item \folder{parallel_state.py}：并行通信组和状态；
  \item \folder{utils.py}：分布式公共工具。
\end{itemize}

它服务于 Tensor Parallel、Pipeline Parallel、Data Parallel、Expert Parallel 和 Context Parallel 等并行模式。

\subsection{\folder{models/}：Ascend 特有模型适配}

这里存放必须在插件侧单独实现或适配的模型，包括 DeepSeek V4、DeepSeek V4 MTP 和 Llama Eagle3 Vwn。注册入口是 \folder{models/__init__.py}。

大多数常规模型仍复用上游 vLLM 模型实现，只把底层执行和算子替换为 Ascend 版本。

\subsection{\folder{patch/}：上游 vLLM 运行时 Patch}

由于本项目是硬件插件，一些上游行为通过运行时 patch 进行适配：

\begin{itemize}
  \item \folder{platform/}：平台、配置、Scheduler、KV Cache、工具解析器等全局 patch；
  \item \folder{worker/}：Worker、模型执行和模型行为相关 patch。
\end{itemize}

Patch 加载过程如下：

\begin{lstlisting}
vLLM loads Ascend plugin
        |
adapt_patch()
        |
import patch/platform or patch/worker
        |
replace or extend upstream behavior
\end{lstlisting}

如果必须只在当前仓库内接入 LoPT，可以利用这一机制包装上游 Renderer/tokenizer，但正式且通用的 tokenization 实现更适合进入上游 vLLM。

\subsection{\folder{quantization/}：量化适配}

该目录提供 Ascend 量化配置和量化方法，包括 FP8、compressed-tensors、ModelSlim、量化类型解析和不同量化 method adapter。它使上游模型加载器能够选择正确的 Ascend 量化执行方式。

\subsection{\folder{sample/}：NPU Token 采样}

负责把模型输出 logits 转换为下一个 token ID，包括 temperature、top-k、top-p、重复惩罚以及推测解码 rejection sampler。

采样处理的是 \code{logits -> next_token_id}，与输入阶段的 \code{text -> token_ids} tokenization 不同。

\subsection{\folder{spec_decode/}：推测解码}

该目录支持 Eagle、Medusa、N-gram、DFlash、DSpark、draft model 和 suffix proposer 等推测解码方案。其目标是先提出多个 draft tokens，再由主模型进行验证，从而提升 Decode 吞吐。

\subsection{\folder{device/}：设备抽象}

负责 Ascend NPU 的设备操作、公共设备工具和 MXFP 类型兼容。

\subsection{\folder{device_allocator/}：显存分配}

负责 CaMem allocator、sleep mode 显存优化以及权重和 KV Cache 内存池管理。

\subsection{\folder{model_executor/}：模型执行辅助能力}

当前主要包含 \folder{offloader/}，用于模型权重 offload、权重预取和 CPU/NPU 之间的模块迁移。

\subsection{\folder{model_loader/}：模型加载扩展}

包含两类额外加载方式：

\begin{itemize}
  \item \folder{netloader/}：通过网络加载或传输权重；
  \item \folder{rfork/}：基于 rfork 的模型和进程加载方案。
\end{itemize}

\subsection{\folder{kv_offload/} 与 \folder{simple_kv_offload/}}

\folder{kv_offload/} 管理 KV Cache 在 CPU 和 NPU 之间的迁移。\folder{simple_kv_offload/} 则适配上游的简单 CPU KV Offload 接口，包含 NPU 内存操作、copy backend 和专用 Worker。两者面向的接入接口和实现路径不同。

\subsection{\folder{eplb/}：专家负载均衡}

EPLB 即 Expert Parallel Load Balancing。该目录负责统计 MoE 专家访问热度、更新专家映射，并在多卡之间平衡专家负载。

\subsection{\folder{lora/}：LoRA 适配}

包含 LoRA 算子、Punica NPU 实现、MoE LoRA 以及 LoRA 张量和索引管理。

\subsection{\folder{profiler/}：性能分析}

集成 Ascend profiler，用于收集 NPU 算子耗时、模型阶段耗时、内存占用和执行性能数据。

\subsection{\folder{_310p/}：Ascend 310P 专用实现}

Ascend 310P 与 910 系列能力存在差异，因此单独维护 Worker、ModelRunner、Attention、采样、MoE、量化和 KV Cache 管理实现。

\subsection{\folder{xlite/}：Xlite 运行时适配}

该目录将 vLLM 模型和权重转换为 Xlite 原生运行时可以识别的形式，提供 \code{XliteWorker}、\code{XliteModelRunner}、模型包装器、KV Cache 注册和 Xlite Graph 执行。这是普通 NPUModelRunner 之外的一条可选执行路径。

\subsection{\folder{_cann_ops_custom/}：自定义算子预留目录}

当前目录中只有 \folder{.gitkeep}，没有实际 Python 逻辑，通常作为构建或安装过程中生成 CANN 自定义算子内容的预留位置。

\section{底层算子：\folder{csrc/}}

\folder{csrc/} 保存 C++、AscendC 和 CANN 扩展代码。

\begin{longtable}{>{\raggedright\arraybackslash}p{0.31\textwidth} >{\raggedright\arraybackslash}p{0.61\textwidth}}
\toprule
\textbf{子目录} & \textbf{功能} \\
\midrule
\endfirsthead
\toprule
\textbf{子目录} & \textbf{功能} \\
\midrule
\endhead
\folder{aclnn_torch_adapter/} & CANN/ACLNN 与 PyTorch Tensor 之间的桥接。 \\
\folder{attention/} & 底层 Attention Kernel。 \\
\folder{batch_matmul_transpose/} & Batch MatMul 和 Transpose 融合算子。 \\
\folder{gmm/} & Grouped MatMul，主要服务于 MoE。 \\
\folder{kernels/} & BGMV、SGMV 等通用 Kernel。 \\
\folder{mc2/} & 计算与通信融合，主要用于 MoE 和分布式执行。 \\
\folder{mla_preprocess/} & MLA 前处理算子。 \\
\folder{moe/} & MoE 专用融合算子。 \\
\folder{common/}、\folder{utils/} & Tensor、CANN、通信和错误处理公共工具。 \\
\folder{cmake/} & 自定义算子的构建逻辑。 \\
\folder{scripts/} & 算子生成和构建辅助脚本。 \\
\folder{third_party/} & 第三方源码或依赖预留位置。 \\
\bottomrule
\end{longtable}

Python 层 \folder{vllm_ascend/ops/} 通常会调用这里编译生成的扩展算子。

\section{测试、示例和开发工具}

\subsection{\folder{tests/}}

\begin{itemize}
  \item \folder{tests/ut/}：单元测试，基本按 \folder{vllm_ascend/} 模块划分；
  \item \folder{tests/e2e/}：端到端推理测试；
  \item \folder{pull_request/}：面向 Pull Request 的基础测试；
  \item \folder{nightly/}：每日完整或硬件相关测试；
  \item \folder{weekly/}：大模型、多节点和低频测试。
\end{itemize}

\subsection{\folder{examples/}}

示例主题包括聊天模板、模型量化、Prefill/Decode 分离、Encoder 解耦、在线 Data Parallel、动态负载均衡、EPLB、rfork 和强化学习集成。示例代码不是在线服务的生产主路径。

\subsection{\folder{benchmarks/}}

用于测量在线吞吐、请求延迟、Serving 性能和算子性能。\folder{tests/*.json} 保存不同性能测试场景的配置。

\subsection{\folder{tools/}}

提供请求发送、Benchmark、自动回归定位、文档生成、代码格式和静态检查工具。代表性文件包括 \folder{send_request.py}、\folder{vllm_bench.py}、\folder{aisbench.py}、\folder{bisect/} 和 \folder{docs_codegen/}。

\section{项目主执行流程}

\begin{lstlisting}
Upstream vLLM CLI / API server
        |
load vllm-ascend plugin
        |
NPUPlatform configuration
        |
NPUWorker.init_device()
        |
load model + allocate KV cache + warm up
        |
SchedulerOutput from upstream vLLM
        |
NPUWorker.execute_model()
        |
NPUModelRunner.execute_model()
        |
Ascend attention / custom operators
        |
logits -> AscendSampler -> next token IDs
        |
upstream output processor / detokenizer
\end{lstlisting}

推荐按照以下顺序阅读源码：

\begin{enumerate}
  \item \folder{setup.py}：了解插件注册；
  \item \folder{vllm_ascend/__init__.py}：了解插件入口和全局 patch；
  \item \folder{vllm_ascend/platform.py}：了解平台和运行配置；
  \item \folder{vllm_ascend/worker/worker.py}：了解 Worker 生命周期；
  \item \folder{vllm_ascend/worker/model_runner_v1.py}：了解每轮推理执行；
  \item \folder{vllm_ascend/attention/} 和 \folder{vllm_ascend/ops/}：了解 NPU 计算实现；
  \item \folder{csrc/}：进一步分析底层融合算子。
\end{enumerate}

\section{与 LoPT 并行 Tokenization 的关系}

当前项目没有正式的生产 tokenization 目录。上游完整请求链路是：

\begin{lstlisting}
HTTP request
    |
chat template / renderer
    |
tokenizer.encode()       <- LoPT candidate location
    |
prompt_token_ids
    |
AsyncLLM / EngineCore / Scheduler
    |
vllm-ascend NPU execution
\end{lstlisting}

如果将 LoPT 原型放入当前仓库，建议采用如下结构：

\begin{lstlisting}
vllm_ascend/
  tokenization/
    __init__.py
    lopt_tokenizer.py
  patch/platform/
    patch_lopt_tokenization.py

tests/ut/
  tokenization/
    test_lopt_tokenizer.py
\end{lstlisting}

其中 \folder{tokenization/} 实现文本分片、overlap、并行调用原 tokenizer 以及边界匹配合并；\folder{patch/platform/} 只负责把它接入上游 vLLM Renderer/tokenizer 调用点。

通常不应为 LoPT 修改以下目录：

\begin{itemize}
  \item \folder{worker/}；
  \item \folder{attention/}；
  \item \folder{ops/}；
  \item \folder{sample/}；
  \item \folder{csrc/}。
\end{itemize}

原因是请求进入这些目录时，原始文本已经被转换为 token IDs，无法再执行基于文本分片边界的 overlap 和匹配合并。

\section{总结}

\folder{vllm-ascend} 的核心职责是完成从上游 Scheduler 生成 token 计算任务，到 Ascend NPU 执行模型并返回新 token IDs 的过程。项目的主运行逻辑集中在 \folder{platform.py}、\folder{worker/}、\folder{attention/}、\folder{ops/} 和 \folder{csrc/}；文档、测试、示例和工具目录则服务于开发、验证和交付。

LoPT 并行 tokenization 位于更靠前的 CPU 请求预处理阶段。当前仓库没有原生 tokenization 主流程，可以通过 patch 机制进行原型接入，但通用实现最终更适合进入上游 vLLM。

\end{document}
